% -*- coding: utf-8 -*-
% !TEX program = xelatex
\documentclass[plain,noanswer,medmath]{../../template/jnuexam}
\usepackage{../../template/kysx}

\begin{document}


\examtitle{1995}{1}


\makepart{填空题}{本题共5个小题，每小题3分，满分15分}


\begin{problem}
$\lim_{x \to 0} (1 + 3x)^{\frac{2}{\sin x}} =$ \fillin{}．
\end{problem}


\begin{problem}
$\frac{\d}{\dx}\int_{x^2}^0x\cos t^2 \dt =$ \fillin{}．
\end{problem}


\begin{problem}
设$(\vec{a}\times \vec{b}) \cdot \vec{c}= 2$,
则$[(\vec{a}+ \vec{b}) \times(\vec{b}+ \vec{c})] \cdot (\vec{c}+ \vec{a}) =$ \fillin{}．
\end{problem}


\begin{problem}
幂级数$\sum_{n = 1}^{\infty}\frac{n}{2^n + (- 3)^n}x^{2n - 1}$的收敛半径$R =$ \fillin{}．
\end{problem}


\begin{problem}
设三阶方阵$A$、$B$满足关系式：$A^{-1} BA = 6A + BA$，且$A = \begin{pmatrix}
  \frac{1}{3} &           0 & 0 \\
            0 & \frac{1}{4} & 0 \\
            0 &           0 & \frac{1}{7}
\end{pmatrix}$，则$B =$ \fillin{}．
\end{problem}


\makepart{选择题}{本题共5小题，每小题3分，满分15分}


\begin{problem}
设有直线$L:\begin{cases}
  x + 3y + 2z + 1 = 0, \\
  2x - y - 10z + 3 = 0
\end{cases}$及平面$\Pi :4x - 2y + z - 3 = 0$，则直线$L$\pickout{}
\begin{abcd}
\item 平行于$\Pi$．   
\item 在$\Pi$上．   
\item 垂直于$\Pi$．   
\item 与$\Pi$斜交．
\end{abcd}
\end{problem}


\begin{problem}
设在$[0,1]$上$f''(x) > 0$，则$f'(0)$,$f'(1)$,$f(1) - f(0)$或$f(0) - f(1)$的大小顺序是\pickout{}
\begin{abcd}
\item $f'(1) > f'(0) > f(1) - f(0)$．       
\item $f'(1) > f(1) - f(0) > f'(0)$．
\item $f(1) - f(0) > f'(1) > f'(0)$．       
\item $f'(1) > f(0) - f(1) > f'(0)$．
\end{abcd}
\end{problem}


\begin{problem}
设$f(x)$可导，$F(x) = f(x)(1 + |\sin x|)$，则$f(0) = 0$是$F(x)$在$x = 0$处可导的\pickout{}
\begin{abcd}
\item 充分必要条件．                
\item 充分条件但非必要条件．
\item 必要条件但非充分条件．        
\item 既非充分条件又非必要条件．
\end{abcd}
\end{problem}


\begin{problem}
设$u_n = (- 1)^n\ln \left(1 + \frac{1}{\sqrt{n}} \right)$，则级数\pickout{}
\begin{abcd}
\item $\sum_{n = 1}^{\infty}u_n$与$\sum_{n = 1}^{\infty}u_n^2$都收敛．
\item $\sum_{n = 1}^{\infty}u_n$与$\sum_{n = 1}^{\infty}u_n^2$都发散．
\item $\sum_{n = 1}^{\infty}u_n$收敛而$\sum_{n = 1}^{\infty}u_n^2$发散．
\item $\sum_{n = 1}^{\infty}u_n$发散而$\sum_{n = 1}^{\infty}u_n^2$收敛．
\end{abcd}
\end{problem}


\begin{problem}
设$A = \begin{pmatrix}
  a_{11} & a_{12} & a_{13} \\
  a_{21} & a_{22} & a_{23} \\
  a_{31} & a_{32} & a_{33}
\end{pmatrix}$,
$$B = \begin{pmatrix}
           a_{21} &          a_{22} & a_{23} \\
           a_{11} &          a_{12} & a_{13} \\
  a_{31} + a_{11} & a_{32} + a_{12} & a_{33} + a_{13}
\end{pmatrix}$$
$P_1 = \begin{pmatrix}
  0 & 1 & 0 \\
  1 & 0 & 0 \\
  0 & 0 & 1
\end{pmatrix}$,
$P_2 = \begin{pmatrix}
  1 & 0 & 0 \\
  0 & 1 & 0 \\
  1 & 0 & 1
\end{pmatrix}$，则必有\pickout{}
\begin{abcd}
\item $AP_1P_2 = B$．  
\item $AP_2P_1 = B$．  
\item $P_1P_2 A = B$．  
\item $P_2P_1 A = B$．
\end{abcd}
\end{problem}


\makepart{计算题}{本题共2小题，每小题5分，满分10分}


\begin{problem}
设$u = f(x,y,z)$, $\varphi (x^2,\e^y,z) = 0$, $y = \sin x$，其中$f$, $\varphi$都具有一阶连续偏导数，且$\frac{\pd \varphi}{\pdz} \ne 0$，求$\frac{\du}{\dx}$．
\end{problem}


\begin{problem}
设函数$f(x)$在区间$[0,1]$上连续，并设$\int_0^1f(x)\dx = A$，求$\int_0^1\dx\int_x^1f(x)f(y)\dy$．
\end{problem}


\makepart{计算题}{本题共2小题，每小题6分，满分12分}


\begin{problem}
计算曲面积分$\iint_{\Sigma}z\dS$，其中$\Sigma$为锥面$z = \sqrt{x^2 + y^2}$
在柱体$x^2 + y^2 \le 2x$内的部分．
\end{problem}


\begin{problem}
将函数$f(x) = x - 1$（$0 \le x \le 2$）展开成周期为$4$的余弦级数．
\end{problem}


\makepart{}{本题满分7分}

\begin{problem}
设曲线$L$位于$xOy$平面的第一象限内，$L$上任一点$M$处的切线与$y$轴总相交，交点记为$A$．已知$\left| \overline{MA} \right| = \left| \overline{OA} \right|$，且$L$过点$\left(\frac{3}{2},\frac{3}{2} \right)$，求$L$的方程．
\end{problem}


\makepart{}{本题满分8分}

\begin{problem}
设函数$Q(x,y)$在$xOy$平面上具有一阶连续偏导数，曲线积分$\int_L 2xy\dx + Q(x,y)\dy$与路径无关，并且对任意$t$恒有
$$\int_{(0,0)}^{(t,1)} 2xy\dx + Q(x,y)\dy = \int_{(0,0)}^{(1,t)} 2xy\dx + Q(x,y)\dy ,$$
求$Q(x,y)$．
\end{problem}


\makepart{}{本题满分8分}

\begin{problem}
假设函数$f(x)$和$g(x)$在$[a,b]$上存在二阶导数，并且
$$g''(x) \ne 0,\quad f(a) = f(b) = g(a) = g(b).$$
试证：
\begin{enumerate}
\item 在开区间$(a,b)$内$g(x) \ne 0$；
\item 在开区间$(a,b)$内至少存在一点$\xi$，使$\frac{f(\xi)}{g(\xi)} = \frac{f''(\xi)}{g''(\xi)}$．
\end{enumerate}
\end{problem}


\makepart{}{本题满分7分}

\begin{problem}
设三阶实对称矩阵$A$的特征值为$\lambda_1 = - 1$,$\lambda_2 = \lambda_3 = 1$，对应于$\lambda_1$的特征向量为$\xi_1 = (0,1,1)^T$，求$A$．
\end{problem}


\makepart{}{本题满分6分}

\begin{problem}
设$A$是$n$阶矩阵，满足$AA^T = E$($E$是$n$阶单位阵，$A^T$是$A$的转置矩阵),$|A| < 0$，求$\left| A + E \right|$．
\end{problem}


\makepart{填空题}{本题共2小题，每小题3分，满分6分}


\begin{problem}
设$X$表示10次独立重复射击命中目标的次数，每次射中目标的概率为0.4，则$X^2$的数学期望$E(X^2) =$ \fillin{}．
\end{problem}


\begin{problem}
设$X$和$Y$为两个随机变量，且
$$P\left\{X \ge 0,Y \ge 0 \right\} = \frac{3}{7},
\qquad P(X \ge 0) = P(Y \ge 0) = \frac{4}{7},$$
则$P\left\{\max(X,Y) \ge 0 \right\} =$ \fillin{}．
\end{problem}


\makepart{}{本题满分6分}

\begin{problem}
设随机变量$X$的概率密度为$f_X(x) = \begin{cases}
 \e^{-x}, & x \ge 0,\\
 0, & x < 0,
\end{cases}$求随机变量$Y = \e^X$的概率密度$f_Y(y)$．
\end{problem}



\end{document}
